% GENERATED FILE. Edit canonical lessons, then run npm run generate:books.
% canonical-source-hash: fnv1a64:333441b4e7077c52
% canonical-lessons: ZH-C429-lutu, ZH-C429-shunlu, ZH-C429-dache-lift, ZH-C429-chuxing, ZH-C429-lufei, ZH-R429-review-425-429, ZH-R429-review-final-429

\chapter{Home Again}
\label{ch:zh-home-again}
\hlchaptermodality{429}

\begin{chapteropening}
\textbf{By the end of this chapter:} \emph{I can talk about the journey, give someone a lift on my way, get a lift, get about and count the costs.}

Complete the last lesson of chapter 429: I can talk about the journey, give someone a lift on my way, get a lift, get about and count the costs.
\end{chapteropening}

\section[lǚtú]{lǚtú — journey}

\label{lesson:ZH-C429-lutu}



\begin{quote}
Before the new one: say the Mandarin for the plan for a trip, then the Mandarin for an official receipt.
\end{quote}

\subsection*{You'll want to know: lǚtú}
\emph{lǚtú} — \textquotedblleft{}journey\textquotedblright{}.

\begin{quote}
\emph{lǚtú hěn cháng wǒmen dōu hěn lèi} — It was a long journey, and we're all tired.
\end{quote}

Its characters come later; for now it is a word for the ear.

\begin{cousinweb}
\emph{lǚtú} is built from \zh{旅} 'to travel' and \emph{tú} 'a road', both inherited from Middle Chinese, \emph{ljoX} and \emph{du} (Guangyun, in Baxter's transcription); the dictionary glosses the first as a body of troops and the second 'a road'.
\end{cousinweb}

\begin{grammarlens}[title={What you've built}]
One more word for getting around.
\end{grammarlens}

\subsection*{Your turn}
\begin{itemize}
\raggedright
  \item \emph{Say it:} \emph{lǚtú}
  \item \emph{Say it:} \emph{lǚtú}, once more
  \item \emph{Say it:} \emph{lǚtú hěn cháng wǒmen dōu hěn lèi} — It was a long journey, and we're all tired.
  \item \emph{Recall:} say the Mandarin for the plan for a trip, then the Mandarin for an official receipt, then say \emph{lǚtú} again
\end{itemize}

\begin{tcolorbox}[breakable,colback=teal!4,colframe=teal!35!black,title={Before you move on}]
What does lǚtú mean? (\textquotedblleft{}Journey\textquotedblright{}.) Say it once more.
\end{tcolorbox}
\section[shùnlù]{shùnlù — on one's way}

\label{lesson:ZH-C429-shunlu}



\begin{quote}
Before the new one: say the Mandarin for an official receipt, then the Mandarin for the journey as a whole.
\end{quote}

\subsection*{You'll want to know: shùnlù}
\emph{shùnlù} — \textquotedblleft{}on one's way\textquotedblright{}.

\begin{quote}
\emph{wǒ zhù zài nǐ jiā fùjìn shùnlù sòng nǐ qù jīchǎng} — I live near you, so I'll drop you at the airport on my way.
\end{quote}

Its characters come later; for now it is a word for the ear.

\begin{cousinweb}
\emph{shùnlù} is built from \emph{shùn} 'to follow' and \zh{路} 'a road', both inherited from Middle Chinese, \emph{zywinH} and \emph{luH} (Guangyun, in Baxter's transcription), which glosses the first 'to follow'.
\end{cousinweb}

\begin{grammarlens}[title={What you've built}]
One more word for getting around.
\end{grammarlens}

\subsection*{Your turn}
\begin{itemize}
\raggedright
  \item \emph{Say it:} \emph{shùnlù}
  \item \emph{Say it:} \emph{shùnlù}, once more
  \item \emph{Say it:} \emph{wǒ zhù zài nǐ jiā fùjìn shùnlù sòng nǐ qù jīchǎng} — I live near you, so I'll drop you at the airport on my way.
  \item \emph{Recall:} say the Mandarin for an official receipt, then the Mandarin for the journey as a whole, then say \emph{shùnlù} again
\end{itemize}

\begin{tcolorbox}[breakable,colback=teal!4,colframe=teal!35!black,title={Before you move on}]
What does shùnlù mean? (\textquotedblleft{}On one's way\textquotedblright{}.) Say it once more.
\end{tcolorbox}
\section[dāchē]{dāchē — to get a lift}

\label{lesson:ZH-C429-dache-lift}



\begin{quote}
Before the new one: say the Mandarin for the journey as a whole, then the Mandarin for on one's way.
\end{quote}

\subsection*{You'll want to know: dāchē}
\emph{dāchē} — \textquotedblleft{}to get a lift\textquotedblright{}.

\begin{quote}
\emph{wǒmen dāchē huí jiā} — We got a lift home.
\end{quote}

Its characters come later; for now it is a word for the ear.

\begin{cousinweb}
\emph{dāchē} is built from \emph{dā} 'to join on' and \zh{车} 'vehicle', both inherited from Middle Chinese, \emph{thap} and \emph{tsyhae} (Guangyun, in Baxter's transcription); the dictionary glosses the first only as a word for touching, so the sense 'to ride along' is not given there.
\end{cousinweb}

\begin{grammarlens}[title={What you've built}]
One more word for getting around.
\end{grammarlens}

\subsection*{Your turn}
\begin{itemize}
\raggedright
  \item \emph{Say it:} \emph{dāchē}
  \item \emph{Say it:} \emph{dāchē}, once more
  \item \emph{Say it:} \emph{wǒmen dāchē huí jiā} — We got a lift home.
  \item \emph{Recall:} say the Mandarin for the journey as a whole, then the Mandarin for on one's way, then say \emph{dāchē} again
\end{itemize}

\begin{tcolorbox}[breakable,colback=teal!4,colframe=teal!35!black,title={Before you move on}]
What does dāchē mean? (\textquotedblleft{}To get a lift\textquotedblright{}.) Say it once more.
\end{tcolorbox}
\section[chūxíng]{chūxíng — to get about; to go out on a journey}

\label{lesson:ZH-C429-chuxing}



\begin{quote}
Before the new one: say the Mandarin for on one's way, then the Mandarin for to get a lift.
\end{quote}

\subsection*{You'll want to know: chūxíng}
\emph{chūxíng} — \textquotedblleft{}to get about; to go out on a journey\textquotedblright{}.

\begin{quote}
\emph{xiàyǔ de shíhou chūxíng hěn bù fāngbiàn} — When it rains, getting about is a real nuisance.
\end{quote}

Its characters come later; for now it is a word for the ear.

\begin{cousinweb}
\emph{chūxíng} is built from \zh{出} 'to go out' and \emph{xíng} 'to go', both inherited from Middle Chinese, \emph{tsyhwit} and \emph{haeng} (Guangyun, in Baxter's transcription).
\end{cousinweb}

\begin{grammarlens}[title={What you've built}]
One more word for getting around.
\end{grammarlens}

\subsection*{Your turn}
\begin{itemize}
\raggedright
  \item \emph{Say it:} \emph{chūxíng}
  \item \emph{Say it:} \emph{chūxíng}, once more
  \item \emph{Say it:} \emph{xiàyǔ de shíhou chūxíng hěn bù fāngbiàn} — When it rains, getting about is a real nuisance.
  \item \emph{Recall:} say the Mandarin for on one's way, then the Mandarin for to get a lift, then say \emph{chūxíng} again
\end{itemize}

\begin{tcolorbox}[breakable,colback=teal!4,colframe=teal!35!black,title={Before you move on}]
What does chūxíng mean? (\textquotedblleft{}To get about; to go out on a journey\textquotedblright{}.) Say it once more.
\end{tcolorbox}
\section[lùfèi]{lùfèi — travel costs; the cost of getting there}

\label{lesson:ZH-C429-lufei}



\begin{quote}
Before the new one: say the Mandarin for to get a lift, then the Mandarin for to go out on a journey.
\end{quote}

\subsection*{You'll want to know: lùfèi}
\emph{lùfèi} — \textquotedblleft{}travel costs; the cost of getting there\textquotedblright{}.

\begin{quote}
\emph{zuò huǒchē qù Běijīng lùfèi shì duōshao} — How much does it cost to get to Beijing by train?
\end{quote}

Its characters come later; for now it is a word for the ear.

\begin{cousinweb}
\emph{lùfèi} is built from \zh{路} 'a road' and \emph{fèi} 'a cost', both inherited from Middle Chinese, \emph{luH} and \emph{phj+jH} (Guangyun, in Baxter's transcription).
\end{cousinweb}

\begin{grammarlens}[title={What you've built}]
One more word for getting around.
\end{grammarlens}

\subsection*{Your turn}
\begin{itemize}
\raggedright
  \item \emph{Say it:} \emph{lùfèi}
  \item \emph{Say it:} \emph{lùfèi}, once more
  \item \emph{Say it:} \emph{zuò huǒchē qù Běijīng lùfèi shì duōshao} — How much does it cost to get to Beijing by train?
  \item \emph{Recall:} say the Mandarin for to get a lift, then the Mandarin for to go out on a journey, then say \emph{lùfèi} again
\end{itemize}

\begin{tcolorbox}[breakable,colback=teal!4,colframe=teal!35!black,title={Before you move on}]
What does lùfèi mean? (\textquotedblleft{}Travel costs; the cost of getting there\textquotedblright{}.) Say it once more.
\end{tcolorbox}
\section[(dialogue)]{Looking back: from The Hotel to Home Again}

\label{lesson:ZH-R429-review-425-429}



\begin{quote}
Say the Mandarin for a place to stay and a big hotel.
\end{quote}

\subsection*{Your turn}
\begin{itemize}
\raggedright
  \item \emph{Say it:} a place to stay, a big hotel, a hotel guest room, a room key card, the hotel lobby
  \item \emph{Say it:} a room with a double bed, fully booked, a small inn, the hotel star rating, a power socket
  \item \emph{Say it:} right next door, to come and go, a buffet, a bed sheet, to be bound for a place
  \item \emph{Say it:} to lead to a place, to set out for a place, the journey back, the plan for a trip, an official receipt
  \item \emph{Say it:} the journey as a whole, on one's way, to get a lift, to go out on a journey, the cost of getting there
  \item \emph{Read it:} \zh{酒} \zh{往} — each one aloud, then a word that holds it
\end{itemize}

\begin{tcolorbox}[breakable,colback=teal!4,colframe=teal!35!black,title={Before you move on}]
A place to stay? (\emph{zhùsù}.) A big hotel? (\textbf{\zh{酒店}}, \emph{jiǔdiàn}.) A hotel guest room? (\textbf{\zh{客房}}, \emph{kèfáng}.) A room key card? (\emph{fángkǎ}.) The hotel lobby? (\emph{dàtáng}.) A room with a double bed? (\emph{dàchuángfáng}.) Fully booked? (\emph{kèmǎn}.) A small inn? (\textbf{\zh{旅店}}, \emph{lǚdiàn}.) The hotel star rating? (\emph{xīngjí}.) A power socket? (\emph{chāzuò}.) Right next door? (\emph{gébì}.) To come and go? (\textbf{\zh{来往}}, \emph{láiwǎng}.) A buffet? (\emph{zìzhùcān}.) A bed sheet? (\emph{chuángdān}.) To be bound for a place? (\textbf{\zh{开往}}, \emph{kāiwǎng}.) To lead to a place? (\textbf{\zh{通往}}, \emph{tōngwǎng}.) To set out for a place? (\textbf{\zh{前往}}, \emph{qiánwǎng}.) The journey back? (\emph{huíchéng}.) The plan for a trip? (\emph{xíngchéng}.) An official receipt? (\textbf{\zh{发票}}, \emph{fāpiào}.) The journey as a whole? (\emph{lǚtú}.) On one's way? (\emph{shùnlù}.) To get a lift? (\emph{dāchē}.) To go out on a journey? (\emph{chūxíng}.) The cost of getting there? (\emph{lùfèi}.)
\end{tcolorbox}
\section[(dialogue)]{Home Again — once more}

\label{lesson:ZH-R429-review-final-429}



\begin{quote}
Say the Mandarin for the journey as a whole and on one's way.
\end{quote}

\subsection*{Your turn}
\begin{itemize}
\raggedright
  \item \emph{Say it:} the journey as a whole, on one's way, to get a lift, to go out on a journey, the cost of getting there
\end{itemize}

\begin{tcolorbox}[breakable,colback=teal!4,colframe=teal!35!black,title={Before you move on}]
The journey as a whole? (\emph{lǚtú}.) On one's way? (\emph{shùnlù}.) To get a lift? (\emph{dāchē}.) To go out on a journey? (\emph{chūxíng}.) The cost of getting there? (\emph{lùfèi}.)
\end{tcolorbox}
